% ============================================================================
%  Relatorio Tecnico - Trabalho 3 de Programacao Paralela (MPI divisao e conquista)
%  Formato exigido:
%    pag. 1  -> coluna DUPLA, margens 1cm, fonte 10, somente texto
%    pag. 2+ -> coluna SIMPLES: tabelas + graficos (anexo A)
%    pag. 3+ -> coluna SIMPLES: codigo fonte
%  Compilar a partir de relatorio/ (usa ../resultados/ e ../ para os fontes):
%    tectonic relatorio.tex    ou    pdflatex relatorio.tex
%  \pend{...} marca em VERMELHO o que falta: numeros do job-extra.sh e as
%  secoes que o grupo escreve (6 a 9 e 11; ver comentarios em cada uma).
%  Antes de entregar: grep pend relatorio.tex deve voltar so a definicao.
%  Os numeros saem de resultados/tabelas.txt (python3 analisar.py).
% ============================================================================
\documentclass[10pt,a4paper,twocolumn]{article}

\usepackage{iftex}
\ifPDFTeX
  \usepackage[utf8]{inputenc}
  \usepackage[T1]{fontenc}
\fi
\usepackage[brazil]{babel}
\usepackage[margin=1cm]{geometry}
\usepackage{graphicx}
\usepackage{booktabs}
\usepackage{listings}
\usepackage{xcolor}
\usepackage{catchfile}
\usepackage{amsmath}
\usepackage{tikz}
\usetikzlibrary{arrows.meta}

\setlength{\parindent}{0pt}
\setlength{\parskip}{2pt}
\setlength{\columnsep}{6mm}
\pagestyle{empty}
% Sem esticar o espaco vertical para encher a pagina (deixava o anexo com um
% buraco no topo) e com titulos de secao do tamanho do texto: a pag. 1 tem de
% caber numa pagina so.
\raggedbottom
\makeatletter
\renewcommand\section{\@startsection{section}{1}{0pt}%
  {0.8ex plus .2ex minus .2ex}{0.3ex}{\normalsize\bfseries}}
\makeatother

\newcommand{\pend}[1]{\textcolor{red}{\textbf{[#1]}}}
\newcommand{\Sdiv}{S_{\mathrm{div}}}
\newcommand{\Ediv}{E_{\mathrm{div}}}

% Inclui o grafico se ja existir; senao deixa um aviso no lugar.
\newcommand{\grafico}[2][\linewidth]{\IfFileExists{../resultados/#2.pdf}%
  {\includegraphics[width=#1]{../resultados/#2.pdf}}%
  {\pend{falta \detokenize{#2}.pdf: rodar python3 graficos.py}}}

% Linhas de tabela geradas por analisar.py (resultados/tab_*.tex), carregadas
% em macros ANTES das tabelas: um \input direto dentro de um tabular quebra o
% \bottomrule ("Misplaced \noalign").
\newcommand{\carrega}[1]{\IfFileExists{../resultados/#1.tex}%
  {\expandafter\CatchFileDef\csname #1\endcsname{../resultados/#1.tex}{}}{}}
\carrega{tab_escala}\carrega{tab_2k}\carrega{tab_muitos}\carrega{tab_ht}\carrega{tab_carga}
\carrega{tab_espera}\carrega{tab_fraca}\carrega{tab_comparacao}
\newcommand{\linhas}[2]{\ifcsname #1\endcsname\csname #1\endcsname\else #2\fi}

\lstdefinestyle{codigoC}{
  language=C, basicstyle=\ttfamily\scriptsize, numbers=left,
  numberstyle=\tiny\color{gray}, keywordstyle=\color{blue!70!black},
  commentstyle=\itshape\color{green!40!black}, stringstyle=\color{red!60!black},
  showstringspaces=false, breaklines=true, tabsize=4,
  columns=fullflexible, lineskip=-1pt,
  literate={á}{{\'a}}1 {ã}{{\~a}}1 {ç}{{\c{c}}}1 {ê}{{\^e}}1 {í}{{\'i}}1
           {ó}{{\'o}}1 {ú}{{\'u}}1 {é}{{\'e}}1 {õ}{{\~o}}1 {à}{{\`a}}1
           {×}{{$\times$}}1 {→}{{$\rightarrow$}}1 {−}{{$-$}}1,
  morekeywords={MPI_Init,MPI_Send,MPI_Recv,MPI_Probe,
                MPI_Get_count,MPI_Status,MPI_Request,MPI_Gather,MPI_Barrier,MPI_Wtime}
}

\begin{document}

% ============================ CABECALHO =====================================
\twocolumn[
{\centering
  {\large\bfseries Bubble Sort com MPI no Modelo Divisão e Conquista}\\[2pt]
  {\small Trabalho 3 --- Programação Paralela --- Grupo 20 (conta LAD \texttt{cp32020})}\\[1pt]
  {\small Gabriel Coelho \quad Gabriel Cherubini}\\[4pt]
}]

% ============================ TEXTO (pag. 1) =================================
% Titulos com as palavras da rubrica. Secoes 1-5 e 10: redigidas com o Claude
% Code (ver secao 10). Secoes 6-9 e 11: escritas pelo grupo; cada uma traz,
% em comentario, o item da rubrica que atende e onde estao os numeros.
\section{Problema}
Ordenar $N$ inteiros com bubble sort (BS) em MPI, por divisão e conquista
numa árvore binária de processos; entrada no pior caso ($N, \dots, 1$),
$N \in \{10^4, 10^5, 10^6\}$, e a raiz confere o resultado. Como o BS é
$\Theta(n^2)$, $L$ pedaços de $N/L$ somam $N^2/L$: dividir já corta
trabalho, e $S > p$ é esperado.

\section{Implementação: mensagens e condição de parada}
Um fonte (\texttt{bs.c}) gera o sequencial e o paralelo
(\texttt{-DPARALELO}); \texttt{-v} escolhe a árvore. Ambas usam o mesmo BS e
a mesma intercalação e trocam $2(p{-}1)$ mensagens bloqueantes: cada rank,
exceto a raiz, recebe um pedaço (tamanho via \texttt{MPI\_Probe} +
\texttt{MPI\_Get\_count}) e o devolve ordenado. Não há mensagem de término:
cada rank sai após devolver o seu, e o programa acaba na última intercalação
da raiz (Fig.~2).

\textbf{Árvore clássica (sem balanceamento).} Filhos de $r$: $2r{+}1$ e
$2r{+}2$. \emph{Parada}: com $\mathit{tam} \le \delta$ o rank ordena e
devolve ao pai; senão envia as metades aos filhos, recebe-as de volta,
intercala e devolve. Com $p = 2^k{-}1$ e $\delta = \lceil N/L \rceil$,
$L = (p{+}1)/2$ folhas, a divisão para exatamente no último nível. Os nós
internos só dividem e intercalam ($O(n)$): ficam quase sempre bloqueados.

\textbf{Árvore do artigo (com balanceamento).} Radenski (PDPTA'11). O rank
$r$ entra no nível $\ell = \lfloor \log_2 r \rfloor + 1$ (raiz: 0) e, a cada
nível, fica com a metade inicial e envia a final ao ajudante $r + 2^{\ell}$,
que a espera com \texttt{MPI\_Probe(MPI\_ANY\_SOURCE)} e adota o remetente
como pai. \emph{Parada}: se $r + 2^{\ell} \ge p$, ordena o que tem; na volta,
recebe de cada ajudante (do último ao primeiro), intercala e devolve ao pai.
Todos ordenam: com $p = 2^k$, pedaços iguais; com $2^k{-}1$, o rank
$2^{k-1}{-}1$ fica com o dobro (4$\times$ o trabalho). Usamos o BS (o
original usa mergesort) e \texttt{MPI\_Send} em vez de \texttt{MPI\_Isend}: o
Open MPI só transfere mensagens grandes quando o remetente volta ao MPI, e pai
e ajudante ordenavam em série.

\section{Metodologia}
Atlântica, 2 nós exclusivos (2$\times$ Xeon E5520: 8 núcleos físicos / 16
hardware threads por nó), gcc 9.4.0 \texttt{-O2}, Open MPI 4.1.1. Ranks em
round-robin: até 8 por nó, um por núcleo; até 16, um por hardware thread
(HT); acima, sem fixação e com \texttt{mpi\_yield\_when\_idle}. Tempo na
raiz, do \texttt{MPI\_Barrier} à última intercalação; mediana de 3. Os
16,6\,min do sequencial de $10^6$, contra 74\,min no enunciado, vêm do
\texttt{-O2} (sem ele, 4,3$\times$ mais lento num teste local).
$S = T_{seq}/T_p$; $E = S/p$ com todos os processos ($S/32$ acima de
32); e $\Sdiv = \sum T_{seq}(\text{pedaço}) / T_p$, o speed-up sobre o
sequencial que ordena \emph{os mesmos pedaços}, sem o ganho de dividir (ideal
$p$; $\Ediv = \Sdiv/p$).

\section{Escalabilidade forte}
O sequencial leva 0,098\,s, 10,0\,s e 993,6\,s. Na clássica com $L$ folhas o
ideal é $S = L^2$ (um $L$ do paralelismo, outro da divisão); com $10^6$ a
medida chega perto até $p = 15$ ($S = 3{,}98$; 15,5; 60,3 para $L^2 = 4$, 16,
64; Tab.~1, Fig.~1). Descontada a divisão, \pend{$\Sdiv$ da clássica $p=15$
e do artigo $p=16$}. Com $10^4$ o custo fixo ($\approx$10\,ms) domina: $E$
cai de 1,17 ($p = 3$) a 0,32 ($p = 31$).

\section{Desbalanceamento da árvore clássica}
Com $p = 15$ e $10^6$, os 7 nós internos ficam 0\,\% ocupados (esperam
17--18\,s, intercalam em 5\,ms) e as 8 folhas, 92\,\%, com pedaços iguais
(máx./média do BS 1,08; Tab.~2, Fig.~3): \pend{$\Ediv$ da clássica $p=15$}.
O problema não é a divisão, e sim os processos parados: só $(p{+}1)/2$
ordenam. Saídas: mais processos que núcleos (Seç.~6) ou internos que também
ordenam (Seç.~7).

% ---------------------------------------------------------------------------
% SECOES 6 A 9 E 11: o grupo escreve (ver "Uso de IA"). Espaco: o que sobrar
% da pagina 1 (compilar e conferir que o Anexo A ainda comeca na pag. 2).
% ---------------------------------------------------------------------------
\section{Mitigação sem alterar o código: aumento do número de processos}
% Rubrica: "como podem ser mitigados ... sem alteração no código" e "desempenho
% ... com balanceamento com aumento do número de processos".
% Responder: o que muda de 15 para 31 processos na classica (16 folhas nos 16
% nucleos fisicos)? Quanto disso e HT (Tab. 4)? Por que a espera dos nos
% internos importa: polling (padrao do Open MPI) x yield, com e sem HT
% (Tab. 5, Fig. 4: tempo total e BS medio de uma folha). Comparar com o
% artigo p=16, que tem os mesmos 16 pedacos. E_div de cada caso (Tab. 8).
\pend{Seção 6, a escrever pelo grupo}

\section{Mitigação com alteração do código: árvore do artigo}
% Rubrica: "como podem ser mitigados ... com alteração no código (versão do
% artigo)" e o desempenho dessa versão.
% Responder: com 2^k todos ordenam pedacos iguais -> S ~ p^2, E_div perto de 1
% (Tab. 3, Tab. 8). Com 2^k - 1 (15: rank 7 com 125.000) o caminho critico e
% o pedaco dobrado: p=15 empata com a classica (16,3 x 16,5 s); p=7 ~ p=4
% (Tab. 1). Carga por nivel do artigo p=15 e p=16 (Tab. 2, Fig. 3).
\pend{Seção 7, a escrever pelo grupo}

\section{63 a 1023 processos}
% Rubrica: "executou as duas versoes com um numero grande de processos (63,
% 127, 255, 512, 1023) e explicou o que aconteceu com o desempenho".
% Responder: o otimo cresce com N (Tab. 6, Fig. 5); o que cresce junto com p
% (lancamento, niveis de comunicacao, disputa por CPU: ~512 processos por no
% em 16 threads com 1023); 63..1023 sao 2^k - 1, entao o artigo sempre tem um
% rank com o pedaco dobrado, e com 64..1024 (2^k) isso some; a partir de 63
% os ranks rodam sem fixacao e com yield (Metodologia).
\pend{Seção 8, a escrever pelo grupo}

\section{Escalabilidade forte × fraca}
% Pedido do grupo: comparar as duas.
% Forte = N fixo, p cresce (Seç. 4): o trabalho por processo cai ~ 1/L^2 e os
% custos fixos dominam; por isso o otimo de p cresce com N (Seç. 8).
% Fraca = n0 = 62.500 por processo ordenador, N = n0 x ordenadores (Tab. 7,
% Fig. 6): o ideal e tempo constante. O que cresce com N: intercalacoes O(N)
% na raiz, envio do vetor, e acima de 16 ordenadores o HT e a espera dos nos
% internos. Por que o speed-up escalado de Gustafson nao vale direto com
% custo quadratico (o sequencial de N cresce com N^2, nao com N).
\pend{Seção 9, a escrever pelo grupo}

\section{Uso de IA}
% CONFERIR com o grupo antes de entregar: tem de descrever exatamente o que
% foi feito pela ferramenta e o que foi feito por cada integrante.
Usamos o \textbf{Claude Code} (Anthropic, modelo Claude Opus 5.5).
\emph{Código}: escreveu o \texttt{bs.c} a partir do pseudocódigo e das
rotinas dadas, a instrumentação de carga e os scripts de medição, análise e
gráficos, e propôs os experimentos complementares (escala fraca, polling
$\times$ yield, $\Sdiv$), aprovados pelo grupo. \emph{Relatório}: redigiu as
Seç.~1--5 e o Anexo A e revisou os números das Seç.~6--9 e 11, escritas pelo
grupo, que também rodou os experimentos no cluster e conferiu resultados e
texto.

\section{Conclusão}
% Rubrica: "avaliou como foi o desempenho da versao sem balanceamento, com
% balanceamento com aumento do numero de processos e com alteracao do codigo".
% Fechar com a comparacao das tres abordagens (Tab. 8).
\pend{Conclusão, a escrever pelo grupo}

% ============================ ANEXO (pag. 2) ================================
\onecolumn
\section*{\large Anexo A --- Tabelas e gráficos}
{\small
\begin{center}
\textbf{Tabela 1.} Escalabilidade forte: tempo (s), speed-up $S$, eficiência
$E = S/p$ e speed-up sobre a mesma divisão ($\Sdiv$, ideal $= p$; $\Ediv = \Sdiv/p$).\\[2pt]
\begin{tabular}{rr rrrrr rrrrr}
\toprule
 & & \multicolumn{5}{c}{árvore clássica (sem balanceamento)} & \multicolumn{5}{c}{árvore do artigo (com balanceamento)}\\
\cmidrule(lr){3-7}\cmidrule(lr){8-12}
$N$ & $p$ & tempo & $S$ & $E$ & $\Sdiv$ & $\Ediv$ & tempo & $S$ & $E$ & $\Sdiv$ & $\Ediv$\\
\midrule
\linhas{tab_escala}{\pend{rodar analisar.py} & & & & & & & & & & & \\}
\bottomrule
\end{tabular}

\vspace{4pt}
\grafico[0.78\linewidth]{grafico_escala}\\[2pt]
\textbf{Figura 1.} Escalabilidade forte: speed-up e eficiência; a linha pontilhada é o ideal linear.
\end{center}

\newpage
\begin{center}
% Diagrama das mensagens: setas cheias = pedaco enviado (com o tamanho); cada
% um volta ordenado pelo mesmo caminho, em sentido inverso.
\begin{tikzpicture}[x=0.85cm, y=1.05cm, font=\scriptsize,
    no/.style={circle, draw=black!60, minimum size=6.5mm, inner sep=0pt},
    interno/.style={no, fill=black!10},
    folha/.style={no, fill=blue!30},
    ordena/.style={no, fill=orange!45},
    passa/.style={no, draw=black!35, dashed, text=black!50},
    envia/.style={-{Stealth[length=2mm]}, semithick},
    fica/.style={densely dotted, black!45}]
  % ---------------- arvore classica, p = 7
  \node[anchor=west] at (-0.6, 0.85) {\textbf{Árvore clássica}, $p = 7$, $\delta = \lceil N/4 \rceil$};
  \node[interno] (c0) at (3, 0) {r0};
  \node[interno] (c1) at (1, -1.3) {r1};
  \node[interno] (c2) at (5, -1.3) {r2};
  \foreach \r/\x in {3/0, 4/2, 5/4, 6/6} \node[folha] (c\r) at (\x, -2.6) {r\r};
  \draw[envia] (c0) -- node[above left=-1pt] {1: $N/2$} (c1);
  \draw[envia] (c0) -- node[above right=-1pt] {2: $N/2$} (c2);
  \draw[envia] (c1) -- node[left] {1: $N/4$} (c3);
  \draw[envia] (c1) -- node[right] {2: $N/4$} (c4);
  \draw[envia] (c2) -- node[left] {1: $N/4$} (c5);
  \draw[envia] (c2) -- node[right] {2: $N/4$} (c6);
  \node[align=left, anchor=north west] at (-0.6, -3.1)
    {nó interno ($\mathit{tam} > \delta$): envia à esquerda (1) e à direita (2),\\
     recebe nessa ordem, intercala e devolve ao pai\\
     folha ($\mathit{tam} \le \delta$): ordena e devolve ao pai};
  % ---------------- arvore do artigo, p = 8
  \begin{scope}[xshift=8.6cm]
  \node[anchor=west] at (-0.6, 0.85) {\textbf{Árvore do artigo}, $p = 8$};
  \node[passa] (a00) at (3.5, 0) {r0};
  \node[passa] (a10) at (1.5, -0.87) {r0};  \node[passa] (a11) at (5.5, -0.87) {r1};
  \node[passa] (a20) at (0.5, -1.73) {r0};  \node[passa] (a22) at (2.5, -1.73) {r2};
  \node[passa] (a21) at (4.5, -1.73) {r1};  \node[passa] (a23) at (6.5, -1.73) {r3};
  \foreach \r [count=\i from 0] in {0,4,2,6,1,5,3,7} \node[ordena] (a3\r) at (\i, -2.6) {r\r};
  \draw[fica] (a00) -- (a10);  \draw[envia] (a00) -- node[above right=-1pt] {$N/2$} (a11);
  \draw[fica] (a10) -- (a20);  \draw[envia] (a10) -- node[above right=-2pt] {$N/4$} (a22);
  \draw[fica] (a11) -- (a21);  \draw[envia] (a11) -- node[above right=-2pt] {$N/4$} (a23);
  \foreach \r/\h in {0/4, 2/6, 1/5, 3/7} {
    \draw[fica] (a2\r) -- (a3\r);
    \draw[envia] (a2\r) -- node[right=-1pt, pos=0.45] {$N/8$} (a3\h);
  }
  \foreach \l/\y in {0/0, 1/-0.87, 2/-1.73, 3/-2.6} \node[black!60] at (7.7, \y) {$\ell = \l$};
  \node[align=left, anchor=north west] at (-0.6, -3.1)
    {em cada nível $\ell$: fica com a metade inicial (pontilhado) e envia a final\\
     ao ajudante $r + 2^{\ell}$; parada: $r + 2^{\ell} \ge p$ $\Rightarrow$ ordena;\\
     na volta recebe de cada ajudante (do último ao primeiro) e intercala};
  \end{scope}
\end{tikzpicture}\\[4pt]
\textbf{Figura 2.} Mensagens nas duas árvores. Cada seta cheia é um pedaço
enviado (com seu tamanho), que volta ordenado pelo mesmo caminho em sentido
inverso: $2(p{-}1)$ mensagens. Na clássica só as folhas (azul) ordenam; no
artigo, todos os processos (laranja; os círculos tracejados são os mesmos
processos em níveis anteriores).

\vspace{8pt}
\textbf{Tabela 2.} Carga média por nível, $N=10^6$ ($p=15$ e $16$): pedaço recebido e
ordenado, tempos (s) e ocupação (BS + intercalação sobre o tempo total).\\[2pt]
\begin{tabular}{lrrr rrrrr r}
\toprule
versão & $p$ & nível & ranks & recebido & ordenado & BS & intercala & espera & ocup.\\
\midrule
\linhas{tab_carga}{\pend{pendente} & & & & & & & & & \\}
\bottomrule
\end{tabular}

\vspace{4pt}
\grafico[0.8\linewidth]{grafico_carga}\\[2pt]
\textbf{Figura 3.} Ocupação de cada processo durante a ordenação ($N=10^6$).

\vspace{8pt}
\begin{minipage}[t]{0.40\linewidth}\centering
\textbf{Tabela 3.} Árvore do artigo com $2^k$ processos.\\[2pt]
\begin{tabular}{rrrrr}
\toprule
$N$ & $p$ & tempo & $S$ & $E$\\
\midrule
\linhas{tab_2k}{\pend{pendente} & & & & \\}
\bottomrule
\end{tabular}
\end{minipage}\hfill
\begin{minipage}[t]{0.58\linewidth}\centering
\textbf{Tabela 4.} HT: tempos (s) de 15 e 31 processos; os mesmos 31 com HT e
só nos núcleos físicos.\\[2pt]
\begin{tabular}{rl rrr rrr}
\toprule
$N$ & versão & $T_{15}$ & $T_{31}$ & 15$\to$31 & HT & físicos & ganho\\
\midrule
\linhas{tab_ht}{\pend{pendente} & & & & & & & \\}
\bottomrule
\end{tabular}
\end{minipage}

\vspace{8pt}
\begin{minipage}{\linewidth}\centering
\textbf{Tabela 5.} Espera dos nós internos: clássica $p=31$ com HT ou só nos
núcleos físicos, com polling (padrão do Open MPI) ou yield; o artigo $p=16$
ordena os mesmos 16 pedaços. BS: tempo de um processo ordenador (média e máximo).\\[2pt]
\begin{tabular}{rlr rrrr}
\toprule
$N$ & configuração & $p$ & tempo & BS méd. & BS máx. & $\Sdiv$\\
\midrule
\linhas{tab_espera}{\pend{job-extra.sh} & & & & & & \\}
\bottomrule
\end{tabular}
\end{minipage}

\vspace{6pt}
\grafico[0.8\linewidth]{grafico_espera}\\[2pt]
\textbf{Figura 4.} Espera ativa (polling) $\times$ ceder a CPU (yield), com e sem HT.

\vspace{8pt}
\begin{minipage}{\linewidth}\centering % titulo e tabela na mesma pagina
\textbf{Tabela 6.} Aumento do número de processos (eficiência sobre as 32 hardware threads).
63 a 1023 são $2^k{-}1$; a última coluna é o artigo com $p{+}1 = 2^k$.\\[2pt]
\begin{tabular}{rr rrr rrr rrr}
\toprule
 & & \multicolumn{3}{c}{árvore clássica} & \multicolumn{3}{c}{árvore do artigo} & \multicolumn{3}{c}{artigo, $p{+}1$ processos}\\
\cmidrule(lr){3-5}\cmidrule(lr){6-8}\cmidrule(lr){9-11}
$N$ & $p$ & tempo & $S$ & $S/32$ & tempo & $S$ & $S/32$ & tempo & $S$ & $S/32$\\
\midrule
\linhas{tab_muitos}{\pend{pendente} & & & & & & & & & & \\}
\bottomrule
\end{tabular}
\end{minipage}

\vspace{4pt}
\grafico[0.8\linewidth]{grafico_muitos}\\[2pt]
\textbf{Figura 5.} Speed-up e eficiência de 1 a 1024 processos.

\vspace{8pt}
\begin{minipage}{\linewidth}\centering
\textbf{Tabela 7.} Escalabilidade fraca: $n_0 = 62.500$ elementos por processo
ordenador, $N = n_0 \times$ ordenadores. BS máx.: o processo ordenador mais
lento; resto: distribuir, intercalar e devolver. $E_\text{fraca} = T_{seq}(n_0)/T$.\\[2pt]
\begin{tabular}{lrrr rrrr}
\toprule
versão & $p$ & ordenadores & $N$ & tempo & BS máx. & resto & $E_\text{fraca}$\\
\midrule
\linhas{tab_fraca}{\pend{job-extra.sh} & & & & & & & \\}
\bottomrule
\end{tabular}
\end{minipage}

\vspace{4pt}
\grafico[0.8\linewidth]{grafico_fraca}\\[2pt]
\textbf{Figura 6.} Escalabilidade fraca: o ideal é tempo constante (eficiência 1).

\vspace{8pt}
\begin{minipage}{\linewidth}\centering
\textbf{Tabela 8.} As três abordagens, $N = 10^6$: sem balanceamento, aumento do
número de processos (sem alterar o código) e árvore do artigo (alterando o código).\\[2pt]
\begin{tabular}{llrr rrrr}
\toprule
abordagem & configuração & $p$ & ordenadores & tempo & $S$ & $\Sdiv$ & $\Ediv$\\
\midrule
\linhas{tab_comparacao}{\pend{rodar analisar.py} & & & & & & & \\}
\bottomrule
\end{tabular}
\end{minipage}
\end{center}
}

% ============================ CODIGO (pag. 3+) ==============================
\newpage
\section*{\large Anexo B --- Código fonte}
\subsection*{bs.c}
\lstinputlisting[style=codigoC]{../bs.c}
\subsection*{lancador.sh, medir.sh e job.sh}
\lstinputlisting[style=codigoC,language=bash]{../lancador.sh}
\lstinputlisting[style=codigoC,language=bash]{../medir.sh}
\lstinputlisting[style=codigoC,language=bash]{../job.sh}
\subsection*{extra.sh e job-extra.sh}
\lstinputlisting[style=codigoC,language=bash]{../extra.sh}
\lstinputlisting[style=codigoC,language=bash]{../job-extra.sh}
\subsection*{analisar.py}
\lstinputlisting[style=codigoC,language=Python]{../analisar.py}
\subsection*{graficos.py}
\lstinputlisting[style=codigoC,language=Python]{../graficos.py}

\end{document}
